\subsection{Macaulay 模的强割子集与商}

设 A 为 Noether 环，M 为有限生成 A-模，S 为 A 的一个子集。按照第八章的约定，我们用 SM 表示 M 的子模 \(\sum_{s\in S}sM\)，并用 G 表示 A 的由 S 生成的理想。

引理 2.--- 设 \(\overline{\mathfrak{G}}\) 为 \(\mathfrak{G}\) 在 A/Ann(M) 中的像。我们有

\[
\operatorname{ht} (\overline {{\mathfrak {G}}}) = \operatorname{codim} (\operatorname{Supp} (\mathrm{M} / \mathrm{SM}), \operatorname{Supp} (\mathrm{M})) .
\]

此外，当 SM ≠ M 时，我们有

\[
\operatorname{ht} (\overline {{{\mathfrak {G}}}}) \leqslant \operatorname{Card} (S).
\]

用 a 表示 M 的零化子。由 II, § 4, 第 4 节, 命题 18 的推论，A-模 M/SM 的支撑集是 V(S + a)。它在 Supp(M) 中的余维数因此等于 V(S + a) 在 V(a) 中的余维数，后者又等于 V((S + a)/a) 在 Spec(A/a) 中的余维数，即 G 的高度。

设 SM ≠ M；当 S 无限时不等式 ht(G) ≤ Card(S) 是显然的，而当 S 有限时它由 VIII, § 3, 第 3 节, 命题 4 b) 得出。

定义 2.--- 设 A 为 Noether 环，M 为有限生成 A-模，S 为 A 的一个有限子集。称 S 对 M 是强割的，如果有

\[
\operatorname{Card} (\mathrm{S}) \leqslant \operatorname{codim} (\operatorname{Supp} (\mathrm{M} / \mathrm{SM}), \operatorname{Supp} (\mathrm{M})) .
\]

注.- 1) A 的每个满足 SM = M 的有限子集 S 都对 M 是强割的。当 SM ≠ M 时，由引理 2，为使 S 对 M 是强割的，必须且只需 Card(S) = codim(Supp(M/SM), Supp(M))，或等价地 ht(G) = Card(S)。

\begin{enumerate}
\def\labelenumi{\arabic{enumi})}
\setcounter{enumi}{1}
\tightlist
\item
  若环 A 是局部环且模 M 非零，则 \(m_{A}\) 的每个对 M 强割的子集 S 都对 M 是割的。事实上，由于 A-模 M/SM 非零，我们有
\end{enumerate}

\[
\operatorname{Card} (S) \leqslant \operatorname{codim} (\operatorname{Supp} (M / S M), \operatorname{Supp} (M)) \leqslant \dim (M) - \dim (M / S M)
\]

(VIII, § 1, 第 2 节, 命题 3 a))，由此得到我们的断言。

命题 5.--- 设 A 为 Noether 环，M 为有限生成 A-模，S 为 A 的一个有限子集。下列条件等价：

\begin{enumerate}
\def\labelenumi{(\roman{enumi})}
\item
  A 的子集 S 对 M 是强割的；
\item
  对 Supp(M/SM) 的每个元素 p，典范映射 \(A \to A_{\mathfrak{p}}\) 诱导一个从 S 到 \(\mathfrak{p}A_{\mathfrak{p}}\) 中一个对 \(M_{\mathfrak{p}}\) 割的子集上的双射。
\item
  (i) \(\Rightarrow\) (ii)：设 \(\mathfrak{p} \in \operatorname{Supp}(\mathrm{M} / \mathrm{SM})\)，并设 \(S'\) 为 S 在 \(A_{\mathfrak{p}}\) 中的像。集合 \(S'\) 含于极大理想 \(\mathfrak{p} A_{\mathfrak{p}}\) 中，并且有
\end{enumerate}

\[
\dim (\mathrm{M} _ {\mathfrak {p}} / \mathrm{S} ^ {\prime} \mathrm{M} _ {\mathfrak {p}}) = \operatorname{codim} (\mathrm{V} (\mathfrak {p}), \operatorname{Supp} (\mathrm{M} / \mathrm{SM}))
\]

(VIII, § 1, 第 4 节, 命题 9)。不等式 Card(S) ≤ codim(Supp(M/SM), Supp(M)) 以及 VIII, § 1, 第 2 节, 命题 3 b) 蕴含关系式

\[
\operatorname{Card} (S) + \dim \left(M _ {\mathfrak {p}} / S ^ {\prime} M _ {\mathfrak {p}}\right) \leqslant \operatorname{codim} \left(V (\mathfrak {p}), \operatorname{Supp} (M)\right) = \dim \left(M _ {\mathfrak {p}}\right).
\]

另一方面，由于 \(M_{p}\) 非零，我们有 \(\dim(M_{\mathfrak{p}}) \leqslant \text{Card}(S') + \dim(M_{\mathfrak{p}}/S'M_{\mathfrak{p}})\) (VIII, § 3, 第 2 节, 公式 (8))。于是由不等式 \(\text{Card}(S') \leqslant \text{Card}(S)\) 得到条件 (ii)。

\begin{enumerate}
\def\labelenumi{(\roman{enumi})}
\setcounter{enumi}{1}
\tightlist
\item
  (ii) \(\Rightarrow\) (i)：不妨设 SM \(\neq\) M。若条件 (ii) 成立，则对 Supp(M/SM) 的每个素理想 \(\mathfrak{p}\)
\end{enumerate}

\[
\operatorname{Card} (S) = \operatorname{Card} \left(S ^ {\prime}\right) \leqslant \dim \left(M _ {\mathfrak {p}}\right) = \operatorname{codim} (V (\mathfrak {p}), \operatorname{Supp} (M)),
\]

对它取下确界即得 (i)。

推论.--- 设 A 为 Noether 环，M 为有限生成 A-模。每个 M-正则序列都对 M 是强割的。

设 x 为一个 M-正则序列，J 为它生成的 \(A\) 的理想。对每个素理想 \(\mathfrak{p} \in \operatorname{Supp}(\mathrm{M/JM})\)，x 在 \(A_{\mathfrak{p}}\) 中的像是 \(\mathrm{M}_{\mathfrak{p}}\)-正则的、由 \(\mathfrak{pA}_{\mathfrak{p}}\) 中元素组成的序列，因而是 \(\mathrm{M}_{\mathfrak{p}}\) 的割序列 (VIII, § 3, 第 2 节, 命题 3 的推论)。

命题 6. 设 A 为 Noether 环，M 为有限生成的 Macaulay A-模，S 为 A 的一个对 M 强割的有限子集。则 A-模 M/SM 是 Macaulay 的。

对每个极大理想 \(\mathfrak{m} \in \operatorname{Supp}(\mathrm{M} / \mathrm{SM})\)，S 在 \(A_{\mathfrak{m}}\) 中的像对 \(M_{\mathfrak{m}}\) 是割的 (命题 5)。由于 \(M_{\mathfrak{m}}\) 是 Macaulay \(A_{\mathfrak{m}}\)-模，\((\mathrm{M} / \mathrm{SM})_{\mathfrak{m}}\) 亦然 (命题 4)，由此得到本命题。

定理 2 (Macaulay-Cohen).--- 设 A 为 Noether 环，M 为有限生成 A-模。下列条件等价：

\begin{enumerate}
\def\labelenumi{(\roman{enumi})}
\item
  A-模 M 是 Macaulay 的
\item
  对 A 的每个由 A 中元素的 M-正则序列生成的理想 J，A-模 M/JM 都没有嵌入伴随素理想；
\item
  对 A 的每个对 M 强割的有限子集 S，A-模 M/SM 都没有嵌入伴随素理想。
\end{enumerate}

(i)⇒(iii)：设 S 为 A 的一个对 M 强割的有限子集。A-模 M/SM 是 Macaulay 的 (命题 6)，特别地没有嵌入伴随素理想 (第 2 节, 命题 2, a))。

(iii)⇒(ii)：这由命题 5 的推论得出。

\begin{enumerate}
\def\labelenumi{(\roman{enumi})}
\setcounter{enumi}{1}
\tightlist
\item
  (ii) → (i)：设 \(\mathfrak{p} \in \text{Supp}(M)\)；我们来证明 \(A_{\mathfrak{p}}\) -模 \(M_{\mathfrak{p}}\) 是 Macaulay 的。对整数 dim( \(M_{\mathfrak{p}}\) ) 作归纳。若 dim( \(M_{\mathfrak{p}}\) ) 为零，则 \(M_{\mathfrak{p}}\) 是有限长度的 \(A_{\mathfrak{p}}\) -模，从而是 Macaulay 的 (第 1 节, 例 1)。设 dim( \(M_{\mathfrak{p}}\) ) 非零，即 p 不是 Supp(M) 的极小元 (VIII, § 1, 第 4 节, 命题 9 a))。由于 M 没有嵌入伴随素理想，p 不含于 M 的任何伴随素理想之中，故存在 p 的元素 x，使乘以 x 的映射 \(x_{M}\) 是单射的 (§ 1, 注 2)。于是乘以 x 的映射 \(x_{M_{\mathfrak{p}}}\) 是单射的，且有 dim( \(M_{\mathfrak{p}}/xM_{\mathfrak{p}}\) ) \textless{} dim( \(M_{\mathfrak{p}}\) ) (VIII, § 3, 第 2 节, 命题 3)。对 A-模 M/xM 以及 Supp(M/xM) 的素理想 p 应用归纳假设，\(A_{\mathfrak{p}}\) -模 \(M_{\mathfrak{p}}/xM_{\mathfrak{p}}\) 是 Macaulay 的，这蕴含 \(A_{\mathfrak{p}}\) -模 \(M_{\mathfrak{p}}\) 是 Macaulay 的 (第 3 节, 命题 4)。
\end{enumerate}

